
\subsection{Implications for Aotearoa New Zealand}
Kaik\={o}ura transfers well in both analyses. Its spatial pattern is reproduced with AUC 0.92 by models that saw no New Zealand data (Figure~\ref{fig:kaikoura}). Its landslide total is forecast within 6\% in the nine-event analysis, where its offset is close to the population mean ($b=-0.12$), and within a factor of two (forecast about half the observed total; $b=+0.62$) in the 38-event analysis. This is encouraging for forecasting a future Alpine Fault or Hikurangi earthquake, but it should not be over-read. The level of landsliding of an unseen earthquake remains uncertain by about two orders of magnitude without mapping. Forecast-time covariates narrow this range (for Kaik\={o}ura the 38-event interval shrinks from 0.03--9.4 to 0.04--4.5 times the observed total, still containing it), but they do not tell in advance whether the next New Zealand earthquake will behave like Cinchona or like Denali. The honest zero-shot statement for a New Zealand scenario is the predictive interval, which spans roughly two orders of magnitude for the landslide total. That interval can be narrowed only by evidence from the event itself. For an Alpine Fault scenario \citep{robinson2016} used in response planning \citep{orchiston2018}, we therefore recommend reporting forecasts as intervals conditioned on scenario ShakeMaps, and planning for the rapid acquisition of a few mapped image tiles as part of the response. Section~\ref{sec:res_fewshot} shows that a handful of tiles is enough to reduce the uncertainty substantially.

Three New Zealand-specific factors are likely to affect the offset of a future event. First, antecedent wetness. The Niigata-Chuetsu earthquake followed heavy typhoon rainfall \citep{sassa2005} and has the second-highest offset in our set ($+1.7$), although our data cannot isolate the contribution of rainfall. Climate projections suggest that the most extreme rainfall-triggered landslide events in New Zealand, such as Cyclone Gabrielle in 2023, may become larger \citep{massey2025gabrielle,dreyer2026climate}. A winter Alpine Fault rupture after a wet period would plausibly sit on the positive side of the offset distribution. Second, inventory practice. The GNS Science inventories are mapped from 0.3~m imagery and lidar \citep{massey2020update}, more completely than several inventories in our training set. All else equal, this should raise Kaik\={o}ura's apparent offset, and in the 38-event analysis, which includes many less complete inventories, its offset is indeed above the mean ($+0.62$). Because the offsets of different inventories are measured relative to each other, the offset of a future New Zealand event will depend on how completely it is mapped as well as on how strongly it shakes the landscape. Third, failure modes that our onshore, source-area model does not represent: runout, which recent work shows is poorly captured by empirical relations for coseismic landslides \citep{grant2025runout}; submarine failures in Kaik\={o}ura Canyon \citep{gnesko2026kaikoura}; landslide volumes \citep{massey2020volume,karantanellis2026kaikoura}; and post-seismic reactivation of slow-moving landslides \citep{dahal2026insar}. Site-specific risk decisions, such as those after Cyclone Gabrielle at Muriwai \citep{howard2025muriwai}, need hazard models at a finer level than the regional forecasts presented here.

\subsection{Limitations and outlook}
Several limitations qualify the results. (1) \emph{Inventory heterogeneity.} The nine inventories differ in mapping resolution, minimum mapped size and polygon conventions \citep{tanyas2017}. Reducing polygons to one representative point gives a common unit but cannot equalise completeness. Part of each event offset is therefore observational rather than physical, and the two cannot be separated with the present data. (2) \emph{Ground motion.} ShakeMap PGA and PGV are smoothed estimates whose quality depends on station density; waveform-based or GNSS-augmented shaking fields can be more informative \citep{dahal2024waveform,endo2026realtime}, and the frequency content of ground motion, which matters for physics-based regional assessment \citep{jiang2026frequency}, is not represented by peak values. (3) \emph{Generic Newmark parameters.} Strength and thickness are uniform; regional geotechnical maps or lithology would sharpen the anchor term. (4) \emph{Static DEM.} GLO-30 post-dates the older earthquakes, so the largest landslide scars may slightly alter the terrain predictors in their source areas. (5) \emph{Leakage and dispersion.} The dispersion factor and the alternative random-effect parameters were estimated from out-of-sample scores of the other events, whose models had seen the target event; the sensitivity analysis in Section~\ref{sec:res_fewshot} shows that coverage depends on $\phi$, whereas point accuracy hardly does, so any leakage mainly affects the stated uncertainty. (6) \emph{Simulated rapid mapping.} Tiles were mapped completely and without error. Operational mapping from optical or SAR imagery \citep{williams2018,handwerger2022} and automated detectors \citep{xu2026rapid,fu2026lsdsam,nasios2026multimodal} will miss landslides and add false positives, which an operational version of the update should model as a detection process. (7) \emph{Number of events.} The 38-event extension supports the main conclusions, but the event covariates were evaluated on few events and were partly derived from metadata of variable quality (for example, station counts for older ShakeMap versions). Adding recent inventories, such as Kahramanma\-ra\c{s} 2023 \citep{yilmaz2026kahramanmaras} and Noto 2024 \citep{zhang2025noto}, and standardised descriptors of mapping resolution would make the covariate-conditioned prior more reliable.
